%───────────────────────────────────────────────────────────────────────────────
% METADATOS DEL DOCUMENTO
%───────────────────────────────────────────────────────────────────────────────

% --- Identificación del documento ---
\newcommand{\docTitulo}{La máquina estocástica}
\newcommand{\docSubtitulo}{Procesos aleatorios,\\ redes neuronales\\ y modelos generativos\\ para físicos}
\newcommand{\docAutor}{J.\,A. Chala Casanova}
\newcommand{\docFecha}{\the\year}
\newcommand{\docEstado}{Versión preliminar}

% --- Cita de portada ---
\newcommand{\docCita}{``El aprendizaje automático es física aplicada a altas dimensiones.''}

% --- Páginas preliminares ---
\newcommand{\docDedicatoria}{A quienes hicieron posible este trabajo.}
\newcommand{\docEpigrafe}{}
\newcommand{\docAgradecimientos}{}

% --- Metadatos PDF ---
\newcommand{\docKeywords}{procesos estocásticos, redes neuronales, modelos de difusión, ecuación de Fokker-Planck, score matching, física computacional, aprendizaje automático}
\newcommand{\docSubject}{Procesos estocásticos y aprendizaje automático para físicos}